\section{Conclusion}
\label{sec:conclusion}

This review updated the project's preliminary research against recent literature on FDG-PET radiomics, ADNI processing pipelines, and cross-validation methodology, resolving three of the five open pipeline-design questions in favor of reusing ADNI's existing processed outputs -- the harmonized \texttt{UCBERKELEYFDG\_8mm} FDG-PET pipeline \cite{jagust2024petcore,jagust2010petcore} and the \texttt{UCSFFSX51} FreeSurfer parcellations -- rather than rebuilding coregistration, normalization, and segmentation from raw images. It found no literature-driven reason to deviate from the project's existing choice of PyRadiomics with IBSI-standard feature classes \cite{vangriethuysen2017pyradiomics,zwanenburg2020ibsi}, but it did surface two changes to the project's current plan that were not previously identified: first, that the Interfile and ECAT dynamic-frame combination step needs decay correction rather than a plain raw-count sum, a methodological gap affecting 216 of the cohort's 2{,}160 image series; second, that the planned classifier stage carries a concrete, quantifiable risk of cross-validation leakage from performing feature selection or class-balancing outside a nested cross-validation loop, an error shown elsewhere in the radiomics literature to inflate reported performance by up to 0.15 AUC-ROC \cite{demircioglu2021bias,demircioglu2024oversampling}.

The review also located a close, recently published methodological precedent for this project's overall cohort-construction strategy -- Mak et al.'s 2025 use of the same Amprion SAA biomarker, timing window, and ADNI FDG-PET data \cite{mak2025alphasynuclein} -- which independently confirms that SAA-positive ADNI subjects carry the expected posterior/occipital hypometabolic signal this project's radiomics pipeline is designed to characterize at finer spatial resolution, and a 2026 multiclass FDG-PET classification study \cite{kim2026dlfdgpet} confirming that separating overlapping/mixed-pathology metabolic patterns, directly analogous to this project's SAA-positive-but-not-clinically-diagnosed subjects, is a recognized open problem in the field rather than an artifact specific to this cohort.

Residual open items not resolved by this review -- the concrete PyRadiomics feature-class list to enable by name, the exact \texttt{src/} directory layout, and the specific nested-cross-validation and feature-selection implementation in scikit-learn -- remain implementation tasks for the pipeline-design stage in \texttt{docs/TODO.md}, to be carried out once the decay-correction fix and the ROI-granularity check described in Section~\ref{sec:methodology} are complete.
